\section{Cut ideals and the Betti-number bound}
\label{sec:counting}

We retain the notation of Section~\ref{sec:laurent}, including the assumption
that the original quiver has a loop at every vertex of positive dimension.
All Hilbert series in this section use half cohomological degree: each
variable $x_d$ has degree one.  We first compute the cut quotients and then
identify the right-hand side of Proposition~\ref{prop:filtration-bound} with
the ordinary Poincar\'e polynomial.
We write $v=\mathbf v$, $w=\mathbf w$, and $l=(l_i)_i$ for the dimension
vectors in this section.

For the counting identities, allow $m=(m_{i,s})_{i\in Q_0,\ s\geq1}$ to be any finite
array of nonnegative integers.  Its slots, cut ideal, and permutation group
are defined as before.  Set
\[
 \|m\|_i=\sum_{s\geq1}s m_{i,s},\qquad
 w(m)=\sum_{i,s}w_i m_{i,s},\qquad
 (q)_m=\prod_{i,s}\prod_{j=1}^{m_{i,s}}(1-q^j).
\]
For two such arrays, write
\[
 b(k,m')=\sum_{i,j,s,t}c_{ij}\min(s,t)k_{i,s}m'_{j,t}.
\]
Thus $D_m=\tfrac12 b(m,m)$.  Empty products equal one, and the empty array
has $P_0=\C$, $J_0=0$, and $h_0(q)=1$.

\begin{lemma}\label{lem:finite-cut-quotient}
For every finite array $m$, the algebra $P_m/J_m$ is finite dimensional.
Consequently $h_m(q)$ is a polynomial with nonnegative integer coefficients.
\end{lemma}

\begin{proof}
The assertion for the empty array is immediate.  Suppose that $x$ is a
common zero of the cut polynomials and that one of its coordinates has a
nonzero value $a$.  Take
\[
 S=\{d:x_d=a\}.
\]
Every framing factor in $q_S(x)$ is nonzero, and every difference across the
cut is nonzero because $S$ contains all slots with value $a$.  Thus
$q_S(x)\ne0$, a contradiction.  The common zero set of $J_m$ is therefore
contained in the origin.  By the Nullstellensatz, either $J_m=P_m$ or a
power of each variable belongs to $J_m$.  Both alternatives give a
finite-dimensional quotient.  Taking invariants under the finite group
$\mathcal W_m$ preserves finite dimensionality and the nonnegative grading.
\end{proof}

\begin{proposition}\label{prop:cut-recurrence}
For every finite array $m$,
\begin{equation}\label{eq:cut-recurrence}
 \frac{1}{(q)_m}
 =\sum_{k+m'=m}
   q^{w(k)+b(k,m')}\frac{h_{m'}(q)}{(q)_k}.
\end{equation}
\end{proposition}

\begin{proof}
We realize the cut products as Euler classes of incoming coordinates and
stratify by unreachable slots to obtain both the coefficient-ring
presentation and the recurrence.  Make a directed graph whose vertices
are the slots of $m$.
For each unordered pair of distinct slots $d,e$, put $b_{de}$ arrows from $d$ to $e$
and $b_{de}$ arrows from $e$ to $d$, with each ordered arrow counted once.
Add a fixed framing vertex and $w_{i(d)}$ arrows from it to $d$.
Assign a one-dimensional vector space to every slot and to the framing
vertex.  Let $E_m$ be the space of representations of this graph, and let
\[
 T_m=(\C^\times)^{\Omega_m^+}
\]
act by change of basis at the slots, leaving the framing line fixed.  The
class $x_d\in H_{T_m}^2(\mathrm{pt})$ is the first Chern class of the standard
character at slot $d$.

A slot is \emph{reachable} if there is a directed path from the framing
vertex to it on which every arrow coordinate is nonzero.  Let $U_m$ be the
open subset where all slots are reachable.  We will show that the natural
coefficient map gives
\begin{equation}\label{eq:cut-cohomology}
 H_{T_m}^*(U_m;\C)=P_m/J_m,
\end{equation}
and that this cohomology is concentrated in even degrees.  Reachability is
an open condition: it is a union of conditions requiring the coordinates
on finitely many paths to be nonzero.

We use induction on the number $n$ of slots.  For $n=0$, the representation
space and its reachable open are both a point, and both assertions hold.
For $0\leq r\leq n$, let $U_r\subset E_m$ be the open subset with at most
$r$ unreachable slots.  Thus $U_n=E_m$ and $U_0=U_m$.  For a subset $S$ of
$r\geq1$ slots, let $Z_S\subset U_r$ consist of representations whose
unreachable set is exactly $S$.

Let $L_S\subset E_m$ be the linear subspace on which all framing arrows
into $S$ and all arrows from the complement of $S$ into $S$ vanish.  These
equations force every slot of $S$ to be unreachable.  Conversely, within
$U_r$, unreachability of all $r$ slots in $S$ forces every complementary
slot to be reachable, so these incoming coordinates must vanish.  It follows that
\[
 Z_S=L_S\cap U_r.
\]
In particular, $Z_S$ is smooth and closed in $U_r$, and
$U_r\setminus U_{r-1}$ is the disjoint union of the $Z_S$ with $|S|=r$.

Let $k$ record the types of the slots in $S$, and put $m'=m-k$.  Restricting
to the complement identifies $Z_S$ with a vector bundle over $U_{m'}$.
Its remaining coordinates are the internal arrows on $S$ and the arrows
from $S$ to its complement, all of which are unrestricted.  Contracting
these vector-bundle fibers is equivariant.  The factor of $T_m$ acting at
the slots of $S$ acts trivially on the base, so the inductive hypothesis
gives
\begin{equation}\label{eq:cut-stratum-cohomology}
 H_{T_m}^*(Z_S)
 =H_{T_{m'}}^*(U_{m'})[x_d:d\in S]
 =(P_{m'}/J_{m'})[x_d:d\in S].
\end{equation}

The normal bundle is the representation formed by the incoming coordinates
set to zero.  Its complex rank and equivariant Euler class are
\begin{align}
 c_S&=w(k)+b(k,m'),\label{eq:cut-stratum-codimension}\\
 e_S&=\prod_{d\in S}x_d^{w_{i(d)}}
       \prod_{d\in S,\ e\notin S}(x_d-x_e)^{b_{de}}
       =q_S.\label{eq:cut-stratum-euler}
\end{align}
Only arrows entering $S$ vanish; this accounts for the exponent
$b(k,m')$, even though the graph has arrows in both directions.

Multiplication by $e_S$ on Equation~\eqref{eq:cut-stratum-cohomology} is
injective.  Indeed, view it as a polynomial in the variables indexed by
$S$, over the possibly nonreduced coefficient ring $P_{m'}/J_{m'}$.  Its
term of highest total degree in those variables is the single monomial
\[
 \prod_{d\in S}x_d^{w_{i(d)}+\sum_{e\notin S}b_{de}}
\]
with coefficient one.  Multiplication by this monomial is injective over
any coefficient ring, and all other terms have lower degree in the new
variables.  A highest-degree argument proves the claim, including the
case $e_S=1$.

Write $i_S:Z_S\hookrightarrow U_r$.  The self-intersection identity says
that $i_S^*i_{S,*}$ is multiplication by $e_S$.  Restriction to a different,
disjoint stratum kills the Gysin image.  Therefore the direct sum of the
Gysin maps is injective.  The equivariant Thom exact sequence, applied
successively from $U_n$ down to $U_0$, gives short exact sequences
\begin{equation}\label{eq:cut-thom-sequence}
 0\longrightarrow
 \bigoplus_{|S|=r}H_{T_m}^{*-2c_S}(Z_S)
 \xrightarrow{\ \sum i_{S,*}\ }H_{T_m}^*(U_r)
 \longrightarrow H_{T_m}^*(U_{r-1})\longrightarrow0.
\end{equation}
It also proves at each step that the cohomology is even and that
restriction from $P_m=H_{T_m}^*(E_m)$ is surjective.  These statements start
with $U_n=E_m$; evenness of the strata follows from the induction on $n$.

We also identify the kernel at each step.  By
Equation~\eqref{eq:cut-stratum-cohomology}, every class on $Z_S$ is the
restriction of a polynomial in the slot variables.  The incoming linear
equations cut out $Z_S$ regularly inside $U_r$, so its fundamental class
$i_{S,*}(1)$ is the restriction of the Euler product $q_S$.  The projection
formula now shows that the image of $i_{S,*}$ is the ideal generated by
that class.  Removing all strata therefore gives exactly the ideal
generated by the nonempty cuts, proving Equation~\eqref{eq:cut-cohomology}.

Finally, choose arrow labels compatibly with permutations of equal types.
Then $\mathcal W_m$ acts on the entire construction, permuting the slot
torus and its variables.  Complex orientations give equivariant Thom
classes without permutation signs.  Taking $\mathcal W_m$-invariants is
exact over $\C$.  Subsets $S$ with a fixed count vector $k$ form one orbit,
whose stabilizer is $\mathcal W_k\times\mathcal W_{m'}$.  Their total
invariant cohomology contribution in
Equation~\eqref{eq:cut-thom-sequence} is consequently
\[
 q^{c_S}
 \Hilb\bigl(\C[x_d:d\in S]^{\mathcal W_k};q\bigr)
 \Hilb\bigl((P_{m'}/J_{m'})^{\mathcal W_{m'}};q\bigr)
 =q^{w(k)+b(k,m')}\frac{h_{m'}(q)}{(q)_k}.
\]
There is one such contribution for each count vector $k$; taking invariants
of the sum over its subset orbit supplies no binomial multiplicity.
The all-reachable open supplies the term $k=0$.  Summing all strata and
using $\Hilb(P_m^{\mathcal W_m};q)=1/(q)_m$ proves
Equation~\eqref{eq:cut-recurrence}.
\end{proof}

\begin{corollary}\label{cor:cut-generating-ratio}
In the formal power series ring $\C(q)[[z_i:i\in Q_0]]$,
\begin{equation}\label{eq:cut-generating-ratio}
 \sum_m z^{\|m\|}q^{-D_m-w(m)}h_m(q)
 =
 \frac{\displaystyle\sum_m
       z^{\|m\|}q^{-D_m-w(m)}/(q)_m}
      {\displaystyle\sum_m z^{\|m\|}q^{-D_m}/(q)_m}.
\end{equation}
\end{corollary}

\begin{proof}
The quadratic form satisfies
\[
 D_{k+m'}=D_k+D_{m'}+b(k,m').
\]
This identity includes the diagonal entries of the type matrix: slots in
the two separate blocks are distinct even when their types agree.
Multiply Equation~\eqref{eq:cut-recurrence} by
$z^{\|m\|}q^{-D_m-w(m)}$ and sum over $m$.  Its summand factors as
\[
 \frac{z^{\|k\|}q^{-D_k}}{(q)_k}\,
 z^{\|m'\|}q^{-D_{m'}-w(m')}h_{m'}(q).
\]
For each monomial in the $z_i$, only finitely many arrays occur.  The
resulting convolution identity may therefore be divided by the first
factor's series, whose constant coefficient is one.  This is
Equation~\eqref{eq:cut-generating-ratio}.
\end{proof}

\subsection{Ordinary cohomology and Hausel's formula}

Let $\nu_{\det}=\prod_i\det_{V_i}$ denote determinant stability, and write
$X_{\det}(v,w)$ for the corresponding quiver variety.  This is a regular
choice: its pairing with every nonzero nonnegative dimension vector is
strictly positive.  Here is the corresponding GIT criterion, including
freeness.  Write $I_i:W_i\to V_i$ for the incoming framing maps.  In the
convention that semi-invariants satisfy $f(gx)=\nu_{\det}(g)^a f(x)$,
semistability is cyclicity: the smallest collection of subspaces invariant
under all doubled-quiver arrows and containing every $\im I_i$ is $V$.
If that collection is proper, a one-parameter subgroup with weight zero
on it and weight $-1$ on a chosen complement has a limit on the point and
negative determinant weight.  Every semi-invariant of positive determinant
weight must therefore vanish.  Conversely, at a cyclic point choose a
basis at each vertex from images of framing vectors under arrow paths.
The product of their basis determinants is a nonzero semi-invariant of
weight $\nu_{\det}$.  On the principal open where this product is nonzero,
normalizing these path bases gives a slice meeting each orbit once.  A
semistable orbit limit is cyclic, and one of its such principal opens also
meets the original orbit; the slice then shows that the limit belongs to
that orbit.  Thus these orbits are closed in the semistable locus.
Moreover, a stabilizer fixes all path images of framing vectors, and hence
is the identity.  With the opposite linearization convention the dual
argument gives cocyclicity.  The argument applies to the closed moment
fiber and allows loops throughout.  Put
\[
 P_{X_{\det}(v,w)}(q)
 =\sum_{j\geq0}\dim H^{2j}(X_{\det}(v,w);\C)q^j.
\]
Odd cohomology vanishes by Kirwan surjectivity, so this is the full ordinary
cohomology polynomial in half degrees.  The empty variety has polynomial
zero.

\begin{proposition}\label{prop:hausel-identification}
For a quiver with a loop at every vertex of positive dimension,
\begin{equation}\label{eq:hausel-cut-bound}
 P_{X_{\det}(v,w)}(q)
 =\sum_{\|m\|=v}
   q^{C(v)-D_m+w\cdot(v-l)}h_m(q).
\end{equation}
\end{proposition}

\begin{proof}
We first specify the moment-level comparison.  Give $N$ a Hermitian
metric invariant under $G_c=\prod_iU(V_i)$, and equip $T^*N$ with its flat
hyperk\"ahler structure.  The hyperk\"ahler moment map has one real and one
complex component.  Under the trace identification of the central real
dual with $\mathbb R^{Q_0}$, let $\eta=(1,\ldots,1)$ be the level corresponding to
the determinant character.  The GIT--moment-map correspondence identifies
$X_{\det}(v,w)$ with the compact-group quotient at real level $\eta$ and
complex level zero.  Rotation of the three moment-map components carries
this quotient to the one at real level zero and complex level $c\eta$,
for a common nonzero normalization constant $c$.  At real level zero,
affine Kempf--Ness identifies the latter quotient with
\[
 \mu_{\C}^{-1}(c\eta)/\!/G.
\]
These are identifications of underlying topological spaces; they do not
assert that rotation is holomorphic in the original complex structure.
The comparison uses the direct rotation on the flat representation
\cite[Section~3(D)]{HKLR87} and the topological affine Kempf--Ness
theorem recalled in Section~\ref{sec:preliminaries}.
The quiver GIT description is explained in \cite{Nakajima94}; the
flat-representation argument just given includes loops.
Multiplying the cotangent coordinate by $c^{-1}$
is an algebraic isomorphism from this affine quotient to
\[
 M_\eta(v,w)=\mu_{\C}^{-1}
       ((\id_{V_i})_i)/\!/G,
\]
because the cotangent moment map is linear in that coordinate.  Thus
$M_\eta(v,w)$ and $X_{\det}(v,w)$ have the same ordinary cohomology,
including simultaneous emptiness.

Hausel's formula for $M_\eta(v,w)$ explicitly allows loops and multiple
arrows \cite[Equation~(9)]{Hausel06}.  We record its grading and all factors
needed here.  For partitions $\lambda^i$ and $\lambda^j$, write
\[
 \langle\lambda^i,\lambda^j\rangle
 =\sum_{a,b}\min(\lambda_a^i,\lambda_b^j),\qquad
 E_Q(\lambda)
 =\sum_{a:i\to j\in Q_1}\langle\lambda^i,\lambda^j\rangle
   -\sum_i\langle\lambda^i,\lambda^i\rangle.
\]
Let $m(\lambda)_{i,s}$ be the number of parts of size $s$ in $\lambda^i$,
and set
\[
 \delta(v,w)=\sum_{a:i\to j\in Q_1}v_iv_j
                  -\sum_i v_i^2+w\cdot v.
\]
The nonempty smooth variety has complex dimension $2\delta(v,w)$.
With $q=t^2$, where $t$ is the ordinary cohomological variable in Hausel's
formula, that formula reads
\begin{equation}\label{eq:hausel-ordinary}
 \sum_v q^{-\delta(v,w)}P_{M_\eta(v,w)}(q)z^v
 =
 \frac{\displaystyle\sum_{\lambda}
       z^{|\lambda|}
       q^{-E_Q(\lambda)-\sum_i\langle\lambda^i,(1^{w_i})\rangle}
       /(q)_{m(\lambda)}}
      {\displaystyle\sum_{\lambda}
       z^{|\lambda|}q^{-E_Q(\lambda)}/(q)_{m(\lambda)}}.
\end{equation}
Here $\lambda$ ranges over all multipartitions, including the empty one.
In particular, the left side uses ordinary cohomology.  Its normalization
can also be checked by Poincar\'e duality: if $u$ is the arithmetic
variable, a smooth variety of complex dimension $2\delta$ satisfies
$P_c(u)=u^{2\delta}P(u^{-1})$.  The normalized count $u^{-\delta}P_c(u)$
therefore becomes $q^{-\delta}P(q)$ after setting $u=q^{-1}$, as in
Equation~\eqref{eq:hausel-ordinary}.

Now write $m=m(\lambda)$.  Removing one loop at each vertex cancels the
vertex self-pairing in $E_Q$, and therefore
\[
 E_Q(\lambda)=D_m,
 \qquad
 \sum_i\langle\lambda^i,(1^{w_i})\rangle
 =\sum_i w_i l_i=w(m),
 \qquad
 \delta(v,w)=C(v)+w\cdot v.
\]
Consequently the right side of Equation~\eqref{eq:hausel-ordinary} is
exactly the ratio in Equation~\eqref{eq:cut-generating-ratio}.
Taking its $z^v$ coefficient and multiplying by $q^{C(v)+w\cdot v}$ proves
Equation~\eqref{eq:hausel-cut-bound}.

The right side is a finite polynomial before this identification:
Lemma~\ref{lem:finite-cut-quotient} makes each $h_m$ finite, there are only
finitely many multipartitions of $v$, and the shifts are nonnegative
integers.  Explicitly,
\[
 C(v)-D_m
 =\tfrac12\sum_{i,j,s,t}c_{ij}
          \bigl(st-\min(s,t)\bigr)m_{i,s}m_{j,t}\geq0,
 \qquad
 w\cdot(v-l)=\sum_{i,s}w_i(s-1)m_{i,s}\geq0.
\]
The diagonal entries $c_{ii}$ are even and the off-diagonal terms pair,
which also verifies integrality.
\end{proof}

\begin{remark}\label{rem:counting-empty-cases}
When $v=0$, only the empty multipartition occurs and the polynomial is one.
When $w=0$ and $v\ne0$, the cut consisting of all positive slots has
$q_S=1$, so every summand is zero.  Equivalently, the numerator and
denominator of Equation~\eqref{eq:cut-generating-ratio} coincide.  This
agrees with emptiness at determinant stability: the diagonal gauge scalar
acts trivially on the fields and nontrivially on the linearization.  An
unframed connected component of the supported graph gives the same unit
cut.  Thus neither the recurrence nor the Betti comparison omits these
degenerate cases.
\end{remark}

\begin{corollary}\label{cor:loop-isomorphism}
Suppose that the quiver has a loop at every vertex of positive dimension.
At determinant stability, the coefficient map induces a canonical
isomorphism of graded $R_0$-algebras
\[
 B_0\xrightarrow{\ \sim\ }H^*(X_{\det}(v,w);\C).
\]
\end{corollary}

\begin{proof}
Proposition~\ref{prop:kernel-inclusion} gives the displayed graded
$R_0$-algebra surjection.  Proposition~\ref{prop:filtration-bound} and
Proposition~\ref{prop:hausel-identification} bound the dimension of each
graded piece of its source from above by that of its target.  Surjectivity
gives the reverse inequality.  Its kernel therefore vanishes in every
degree.  Since the bijection is the original coefficient-compatible
algebra map, this is an isomorphism over $R_0$ and an equality of the two
coefficient ideals, with the quotient rings retaining nilpotents.  The same argument
includes the zero ring in the empty case.
\end{proof}
